% !TeX program = lualatex
\documentclass[12pt,oneside]{book}
\usepackage[utf8]{inputenc}
\usepackage[T1]{fontenc}
\usepackage{geometry}
≥ometry{margin=1in}
\usepackage{amsmath,amssymb}
\usepackage{graphicx}
\usepackage{xcolor}
\usepackage[dvipsnames]{xcolor}   % for coloured suits
\usepackage{amsfonts, amssymb}    % for suit symbols
\usepackage{tcolorbox}
\tcbuselibrary{skins, breakable}

% −−−−− Card suit and rank formatting −−−−−
% Define suit colours (hearts/diamonds red, spades/clubs black)
≠wcommand{\cardSuit}[1]{%
  ℹfnum\pdfstrcmp{#1}{♠}=0 \textcolor{black}{♠}%
  \elseℹfnum\pdfstrcmp{#1}{♥}=0 \textcolor{red}{♥}%
  \elseℹfnum\pdfstrcmp{#1}{♦}=0 \textcolor{red}{♦}%
  \elseℹfnum\pdfstrcmp{#1}{♣}=0 \textcolor{black}{♣}%
  \else #1\fi\fi\fi\fi
}

% Command to print a full card: \card{♠}{A} or \card{♠}{10}
≠wcommand{\card}[2]{%
  \tcbox[colback=white, colframe=gray!50, sharp corners, boxsep=1pt, left=2pt, right=2pt, top=1pt, bottom=1pt]%
    {\cardSuit{#1}\ \textbf{#2}}%
}

% Optional: shortcut commands for each suit
≠wcommand{\spadecard}[1]{\card{♠}{#1}}
≠wcommand{ℎeartcard}[1]{\card{♥}{#1}}
≠wcommand{\clubcard}[1]{\card{♣}{#1}}
≠wcommand{⋄card}[1]{\card{♦}{#1}}


% Suit symbols – used only in imported generator files
≠wcommand{♠}{♠}
≠wcommand{♥}{♥}
≠wcommand{♣}{♣}
≠wcommand{♦}{♦}

% Almanac environment
≠wenvironment{almanac}
  {\begin{quote}\smallℹtshape}
  {\end{quote}}

\begin{document}

% −−−−−−−−−−−−−−−−−−−−−−−−−−−−−−−−−−−−−−−−−−−−−−−−−−−−−−−−−−−−−−−−−−−−−−
% FRONT MATTER
% −−−−−−−−−−−−−−−−−−−−−−−−−−−−−−−−−−−−−−−−−−−−−−−−−−−−−−−−−−−−−−−−−−−−−−
\frontmatter

\title{\Huge Fate's Edge\\[0.3\baselineskip]
Łarge Dhahara (The Kingdoms of Brass)}
\author{Viceroy’s Scribe Govind \\ \small (keeper of the iron ledger, Sarvistan weigh‑house)}
\date{→day}
\maketitle

\tableofcontents

% −−−−−−−−−−−−−−−−−−−−−−−−−−−−−−−−−−−−−−−−−−−−−−−−−−−−−−−−−−−−−−−−−−−−−−
% PREFACE – Framing device (native Dhaharan scribe)
% −−−−−−−−−−−−−−−−−−−−−−−−−−−−−−−−−−−−−−−−−−−−−−−−−−−−−−−−−−−−−−−−−−−−−−
\chapter*{Preface: The Iron Ledger}
\addcontentsline{toc}{chapter}{Preface}

\noindent
{Łarge
\textit{The mountain weighs every breath, every ingot, every oath. Pay in iron, pay in silence, or pay in blood. And when the ninth lease cracks, even the elephants will remember who fed them during the famine.}
}

I am a viceroy’s scribe – Dhaharan by birth, Aeler by employment. My grandfather kept the toll accounts at the Himadri passes; my mother balanced the grain quotas for the Brass Coast. I inherited an iron stamp and a healthy suspicion of anyone who does not seal their contracts with wax and a witness.

Dhahara is not a mystery. It is a negotiation – between the lowland kings and the Aeler viceroys, between the hill rajas and the Admiralty, between the elephant and the mahout. The Brass Coast runs on tariffs, not trust. The Nine Kingdomss measure power in elephant herds and marriage alliances. The Himadri foothills measure it in breaths – eight is safe, nine is the mountain’s.

You will find no abstract “debt” here. You will find iron ingots stamped with a viceroy’s mark, pepper auctions where a sneeze can cost a fortune, and a sealed mine where a demon whispers promises it cannot keep. Pay attention to the ninth clause. It is always the one that kills you.

The ledger is open. The scales are balanced. The ninth seal is warm.

\vspace{1em}
\noindent\textit{— Viceroy’s Scribe Govind, Sarvistan weigh‑house, stamped and witnessed}

% −−−−−−−−−−−−−−−−−−−−−−−−−−−−−−−−−−−−−−−−−−−−−−−−−−−−−−−−−−−−−−−−−−−−−−
% HOW TO USE THIS BOOK
% −−−−−−−−−−−−−−−−−−−−−−−−−−−−−−−−−−−−−−−−−−−−−−−−−−−−−−−−−−−−−−−−−−−−−−
\chapter*{How to Use This Book}
\addcontentsline{toc}{chapter}{How to Use This Book}

This book is a companion to the \textit{Fate’s Edge Core Amaranthine Worldbook} and the \textit{System Reference Document (SRD)}. It provides regional content for Dhahara – the Brass Coast, the Nine Kingdomss, and the Himadri Foothills – using the same card‑driven generator system.

Each regional chapter contains:
\begin{itemize}
  ℹtem A short introduction (in the voice of Govind, who knows where the bodies are buried – literally).
  ℹtem A local quote.
  ℹtem An \emph{almanac} – a quick Q&A about daily life in that region.
  ℹtem Then the full generator tables (♠\ places, ♥\ people, ♣\ complications, ♦\ rewards), imported from separate files. To use them, draw cards and consult the appropriate rank as described in the Core book’s \textit{Cartomancer’s Compass}.
\end{itemize}

For the core rules (attributes, skills, actions, magic, summoning, travel, etc.), refer to the \textit{Fate’s Edge SRD or Player/GM Guides}. The generators assume you are familiar with Position, Effect, and Timers.

\vspace{1em}
\noindent\textit{“The road is a ledger – but only if you forget the elephants.”} – Dusana of the Raven Road (apocryphal)

% −−−−−−−−−−−−−−−−−−−−−−−−−−−−−−−−−−−−−−−−−−−−−−−−−−−−−−−−−−−−−−−−−−−−−−
% A NOTE ON DHARANAN CULTURES (sensitivity guidance)
% −−−−−−−−−−−−−−−−−−−−−−−−−−−−−−−−−−−−−−−−−−−−−−−−−−−−−−−−−−−−−−−−−−−−−−
\chapter*{A Note on Dhaharan Cultures}
\addcontentsline{toc}{chapter}{A Note on Dhaharan Cultures}

Dhahara is a subcontinent of immense diversity – coastal spice ports, riverine kingdoms, highland iron mines, and steppe edges. The cultures presented here draw inspiration from South Asian history and geography, but they are not direct analogues. This note offers guidance for respectful portrayal.

\subsection*{Avoiding Monoliths}
The Nine Kingdomss are not a single “India”. They are rival dynasties with distinct customs, languages, and patron deities. A queen of Vajrapur may be a devout follower of Mykkiel; a king of Mahishmati may worship a local buffalo spirit. Show internal variety.

\subsection*{Agency, Not Exoticism}
The Dhaharan rajas are not passive subjects of the Aeler viceroyalty. They negotiate, resist, and exploit the iron trade. The Sidhi diaspora on the Brass Coast are not helpless refugees – they run ropeways, lighthouse courts, and freedman’s seals. Let players see Dhaharan characters as competent, proud, and politically savvy.

\subsection*{Elephants as Partners, Not Props}
War elephants are not dumb beasts. They have names, lineages, and mahouts who speak to them in a language of clicks and whistles. An elephant may remember a kindness for decades – and a betrayal longer.

\subsection*{The Ninth Seal}
The demon in the Himadri mines is not a “dark god” caricature. It is a fragment of Varash the Unmaker, a being of cosmic indifference. It does not demand worship; it offers bargains. The horror is in the fine print.

\subsection*{GM Mantra}
“Portray dignity, complexity, and the weight of a kept promise.”

% −−−−−−−−−−−−−−−−−−−−−−−−−−−−−−−−−−−−−−−−−−−−−−−−−−−−−−−−−−−−−−−−−−−−−−
% REGIONAL CHAPTERS
% −−−−−−−−−−−−−−−−−−−−−−−−−−−−−−−−−−−−−−−−−−−−−−−−−−−−−−−−−−−−−−−−−−−−−−
\mainmatter

\chapter{Gateways & Travel}
\emph{The elephant remembers. The mountain forgets. The demon waits.}

\section*{Travel Modes of Dhahara}

\subsection*{Coast Roads (Brass Coast)}
Tariffs, ropeways, and lighthouse courts. A \textbf{Lighthouse Charter} (Sidhi) or a \textbf{Customs Seal} (Admiralty) buys passage. Never board a ship with a green lantern at the mast – it means the captain is under a death vow.

\subsection*{River Routes (Nine Kingdomss)}
The rivers of the Nine Kingdomss are borders and battlegrounds. A \textbf{Peacock Feather} or an \textbf{Iron Gate Pass} grants audience with a king. Never praise a king’s elephant – compliment the mahout instead.

\subsection*{High Passes (Himadri Foothills)}
Breath‑counting, altitude sickness, and the Ninth Seal. A \textbf{Breath‑Counter} or a \textbf{Toll Waiver} from the Aeler viceroy is required. Never count your breaths aloud in a mine – the mountain listens.

\subsection*{Elephant Caravans}
The walking cities of Ngomebe are distant neighbours, but Dhaharan elephant caravans carry iron and spices along the \textbf{Salt Scar}. A \textbf{Keystone Shard} or a \textbf{Mason‑Kin Hammer} may convince the Stone Men to share their path.

\section*{Gateways to Other Expansions}

\begin{itemize}
  ℹtem \textbf{Amaranthine Core (Kahfagia, Ecktoria, etc.)}: From the Brass Coast, the \textbf{Black River} leads north to Thepyrgos. From the Himadri passes, the \textbf{Aeler Underways} connect to the Spine.
  ℹtem \textbf{Southern Realms (Oshiira, Ashaan, Ngomebe)}: The Brass Coast is a direct sea route to Oshiira and Ashaan. The \textbf{Salt Scar} connects the Nine Kingdomss to Ngomebe’s walking cities.
  ℹtem \textbf{Way of Silk (Fhara, Sidhi, Pereshi, Kuvani, Tulkani)}: The Brass Coast is home to a large Sidhi diaspora; lighthouse charters are accepted as credit. The Himadri passes lead to Pereshi highlands.
  ℹtem \textbf{Eastern Realms (Sihai, Nihori, Ayokha)}: Monsoon winds from the Brass Coast carry ships to Ayokha. The \textbf{Silk Road} east from the Nine Kingdomss leads to Sihai.
\end{itemize}
≠wpage

% −−−−−−−−−−−−−−−−−−−− THE BRASS COAST −−−−−−−−−−−−−−−−−−−−
\chapter{The Brass Coast}
\emph{Shore of Silver and Spice.}

\noindent
Govind says: \textit{“Three powers, one coast, and a tariff schedule that changes with the tide. The Admiralty holds the ports. The Sidhi run the ropeways. The rajas sharpen their swords in the hills. Choose a side – or let the coast choose for you.”}

\noindent
The Brass Coast is Dhahara’s southern rim: a humid, feverish coastline of palm groves, salt pans, and crowded port cities. Three powers bleed into each other here, and none can draw a clean breath. The Kahfagian Admiralty holds the ports and the tariff ledgers. The Sidhi diaspora runs the ropeways, the lighthouse courts, and the freedman’s seal. And the native Dhaharan rajas – the hill lords who have lived here since before any foreign flag flew – grow the pepper, mine the iron, and watch the coast from inland fortresses, waiting. The air smells of pepper, cinnamon, and the low tide of a thousand ships.

\vspace{0.5em}
\noindent\textit{“Never board a ship that has a green lantern at the mast. It means the captain is under a death vow – and you are the payment.”} \\
— Raja Govind of the Kalinga foothills

\vspace{1em}
\begin{almanac}
\textbf{What do people eat for breakfast?} Rice and fermented fish paste – the rice from coastal paddies, the fish caught by lighthouse keepers’ children. The children are paid in copper. The ninth child never takes payment. No one knows why.

\textbf{What does a child learn before they can walk?} To read a lighthouse code. The eight lamp‑flashes: Tariff, Inspection, Quarantine, Smuggler, Pirate, Rescue, Home, and the ninth – the flash that means the collector is coming. The ninth flash is never used. It is a warning.

\textbf{What is the first thing you notice about a stranger?} Their papers. The Admiralty requires a manifest for every shipment. A stranger with perfect papers is a spy. The ninth page of every manifest is blank – that is where the dead travel.

\textbf{What is the most common crime?} Tariff evasion – but evasion is a profession. The tariffs are set by the Admiralty. The evaders pay the Admiralty directly. The payment is called a “consulting fee.” The ninth fee is always paid in salt, not silver.

\textbf{What does no one talk about but everyone knows?} The lighthouse lamps are not fuelled by oil. They are fuelled by the fat of the drowned – whales, sailors, slaves who did not survive the crossing. The fat burns green. The ninth lamp on every tower is never lit. It is reserved for the collector.

\textbf{Where do people go when they die?} To the sea. Weighted with brass and dropped in the deep channel called the Freedman’s Grave. The ninth body dropped each year is never weighted. It always drifts toward the Admiral’s palace.

\textbf{How do you apologize for a serious wrong?} You offer a tariff – a payment of silver, calculated by the wronged party’s estimate of their loss. The court takes a percentage. The ninth tariff ever collected was paid in a memory. The judge still remembers it.

\textbf{What is the most important festival?} The Burning of the Ledgers (new year). Every merchant’s ledger is brought to the lighthouse. The ledgers are read aloud, then burned. The ash falls on the water. The ninth ledger is never burned – it is kept in a brass box at the bottom of the Freedman’s Grave.

\textbf{What is the secret that could destroy a powerful person?} The lighthouse lamps are not all green. One lamp burns blue – the Admiral’s personal beacon. The blue flame is a spirit – a Sidhi oracle bound to the lighthouse nine generations ago. Her ninth prophecy was never recorded. It said: \textit{“The collector will come by sea, and the ledger will burn from the bottom up.”}
\end{almanac}

\subsection*{Integration Notes}
The Brass Coast is a hub for the \textbf{Way of Silk} – Sidhi lighthouse charters are accepted as proof of credit from Fhara to Kuvani. Taharkan coffee arrives on monsoon winds, traded for brass and iron. Ayokhan spice merchants maintain warehouses here. And Sihai silk convoys offload at the same piers, paying tariffs to the Admiralty and bribes to the rajas.

≠wpage
∈put{../generators/brass_coast.tex}
≠wpage

% −−−−−−−−−−−−−−−−−−−− THE NINE KINGDOMS −−−−−−−−−−−−−−−−−−−−
\chapter{The Nine Kingdomss}
\emph{The Peacock Throne and Its Rivals.}

\noindent
Govind sighs: \textit{“The number nine is a lie. There have never been nine kingdoms. The number is a prayer – a hope that the Hollow will not notice that one is always missing. The throne weeps because the missing kingdom is always the one that mattered.”}

\noindent
Between the Satrapies plateau and the Brass Coast lies a fertile, contested lowland: the land of the Nine Kingdomss, the heart of Dhahara. Here, ancient dynasties still claim descent from the sun, and the Peacock Throne of Ahimsa casts its shadow over half a dozen rival courts. Rivers are borders; elephants are war machines; every harvest is a negotiation between kings who have never met but whose grandfathers swore blood oaths. The Covenant has a strong presence in the sacred cities; the Aeler have counting houses in the iron towns; the Sidhi sail up the rivers to trade. But Dhahara belongs to the Dhaharan – to their poets, their temple dancers, their bandit kings, and their queens who rule from behind silk screens.

\vspace{0.5em}
\noindent\textit{“Never praise a king’s elephant. The elephant is sacred; your praise will be taken as an insult to the king’s own mount. Compliment the mahout instead. The mahout will tell the elephant.”} \\
— Temple Dancer Mira

\vspace{1em}
\begin{almanac}
\textbf{What do people eat for breakfast?} Rice porridge with mango pickle – the rice from the royal granaries, the mangoes from the king’s orchards. The ninth grain of rice in every bowl is left uneaten. It is an offering to the missing kingdom.

\textbf{What does a child learn before they can walk?} To recite their genealogy. Nine Kingdomss children learn the names of their ancestors back to the sun. The ninth ancestor is never named. That is the gap where the Hollow sits.

\textbf{What is the first thing you notice about a stranger?} Their caste mark. The nine castes – priest, warrior, merchant, farmer, artisan, servant, outcast, foreigner, and the ninth – the caste that has no name. The ninth caste is the king’s spies. Their faces are always the same. No one can describe them.

\textbf{What is the most common crime?} Tax evasion – a political statement. The king’s tax collectors are also his spies. The ninth tax collector vanished last year. His shadow still walks the counting house.

\textbf{What does no one talk about but everyone knows?} The Peacock Throne is empty. The Maharaja died seven years ago. The courtiers rule in his name. The ninth minister is a eunuch who controls the palace guards. He has no shadow.

\textbf{Where do people go when they die?} To the river. Cremated on the ghats, ashes scattered on the water. The ninth ghat is reserved for the missing kingdom. It has never been used. It is always warm.

\textbf{How do you apologize for a serious wrong?} You offer a dance – a performance at the temple of the sun. If the priest smiles, you are forgiven. If the priest weeps, you are not. If the priest is silent, you are cursed. The ninth dance ever performed was for a debt to the Hollow. The dancer is still dancing.

\textbf{What is the most important festival?} The Coronation of the Monsoon (summer). The king offers a golden elephant to the rain clouds. The clouds have refused for nine years. The ninth golden elephant was hollow. It was filled with the ashes of the ninth king.

\textbf{What is the secret that could destroy a powerful person?} The golden elephant is hollow. It is filled with the ashes of the ninth king – the one who tried to unite the kingdoms and failed. His ghost rides the monsoon winds. He is still trying.
\end{almanac}

\subsection*{Integration Notes}
The Nine Kingdomss are a crossroads for the \textbf{Way of Silk} – Tulkani caravans pass through, trading story‑knots for elephant ivory. Taharkan coffee is a favourite of the Maharaja’s court. Sihai silk is measured against Aeler iron in the counting houses of Vajrapur. And the Kalrantha amulet – a fragment of the demon Varash – is said to be hidden somewhere in the ruins of Rajagriha, sought by thieves and Veiled Sisters alike.

≠wpage
∈put{../generators/nineₖingdoms.tex}
≠wpage

% −−−−−−−−−−−−−−−−−−−− THE HIMADRI FOOTHILLS −−−−−−−−−−−−−−−−−−−−
\chapter{The Himadri Foothills}
\emph{The Roof of the World.}

\noindent
Govind shivers: \textit{“The mountain does not speak in words. It speaks in the silence after a bell, the weight of a frozen lung, the taste of iron on the tongue at dusk. We Aeler do not conquer stone. We listen to it – and we remember what it forgets.”}

\noindent
The Himadri foothills are the highest inhabited place in the known world. Here, the Aelerian mountains scrape the sky, and the air is so thin that each breath is a measured gift. The native Aeler of these peaks – called the High Aeler or the Karak‑Aeler (“Stone‑Breath People”) – are not the mercantile dwarves of the lowland holds. They are severe, cold, and ancient, their faces carved by wind and frost, their voices rarely raised above a whisper. They do not trade in breath‑bonds for profit; they trade in survival. The humans of the Himadri – the High Folk, or Mi‑tsang – are equally remote: a scattering of monastic villages, goat‑herding clans, and hidden fortress‑temples that have never paid tribute to any king. The demon sealed in the Ninth Seal Mine – a fragment of Varash the Unmaker – is not a legend. It is a fact. The seal is cracking. And the mountain is listening.

\vspace{0.5em}
\noindent\textit{“Never count your breaths aloud in a mine. The mountain listens, and if it hears you reach the ninth, it will take one – a breath, a memory, a shadow.”} \\
— Mine‑Breath, who carries her grandfather’s bell

\vspace{1em}
\begin{almanac}
\textbf{What do people eat for breakfast?} Iron rations and tea – the rations are hardtack and dried meat, the tea brewed from leaves that grow only above the snow line. The leaves are harvested by the Vent‑Prior’s apprentices. The lottery is a death sentence.

\textbf{What does a child learn before they can walk?} To read a pressure gauge. The eight pressures: Safe, Warning, Danger, Collapse, Flood, Fire, Seal, and the ninth – the pressure that means the mountain is breathing. The mountain’s breath is the miners’ death.

\textbf{What is the first thing you notice about a stranger?} Their bell. Every miner carries a breath‑counter bell. A bell that rings too fast means running. A silent bell means dead. A bell that rings the ninth note means a demon.

\textbf{What is the most common crime?} Ore smuggling – but smuggling is a perquisite. The mine‑captains take a cut called “the Vent‑Prior’s fee.” The negotiation involves tea and threats.

\textbf{What does no one talk about but everyone knows?} The Ninth Seal is not a seal. It is a door – a gate to a prison that holds something that should not be named. The keepers do not want to know what is inside.

\textbf{Where do people go when they die?} To the mountain. Bodies left in worked‑out galleries. The mountain absorbs them. Its memory grows.

\textbf{How do you apologize for a serious wrong?} You offer a breath – one of your own, counted and witnessed. Sealed in a brass capsule and placed in the Vent‑Prior’s ledger. The ledger is the court. The mountain’s judgment is final.

\textbf{What is the most important festival?} The Counting of the Breath (winter solstice). Every miner’s breath‑counter is read aloud. Those with the lowest counts are promoted. Those with the highest are retired to a sealed gallery. The gallery is a tomb.

\textbf{What is the secret that could destroy a powerful person?} The demon in the Ninth Seal is not a demon. It is an Aeler – the first Vent‑Prior, who volunteered to be sealed with the mountain’s hunger. He is still alive. He has been counting his breaths for nine centuries. He has reached the ninth count. He is still counting.
\end{almanac}

\subsection*{Integration Notes}
The Himadri iron mines feed the forges of the Brass Coast and the walking cities of Ngomebe. Aeler breath‑bonds are recognised by Covenant judges as valid contracts. The Ninth Seal demon is whispered about by Sihai monks (who have their own sealed archives) and by Breath‑less agents (who believe it is a fragment of the same entity bound in the Kalrantha amulet). The Kalrantha amulet itself is said to have been forged in these mountains – a green glass shard that can still be found in the scree if you know where to look.

≠wpage
∈put{../generators/himadri_foothills.tex}
≠wpage

% −−−−−−−−−−−−−−−−−−−−−−−−−−−−−−−−−−−−−−−−−−−−−−−−−−−−−−−−−−−−−−−−−−−−−−
% APPENDICES
% −−−−−−−−−−−−−−−−−−−−−−−−−−−−−−−−−−−−−−−−−−−−−−−−−−−−−−−−−−−−−−−−−−−−−−
\appendix

% −−−−−−−−−−−−−−−−−−−− TRADE & TARIFFS −−−−−−−−−−−−−−−−−−−−
\chapter{Trade & Tariffs of Dhahara}

\subsection*{The Spice Auction (Brass Coast)}
Pepper, cardamom, and cinnamon are sold by candle auction: a one‑inch candle is lit; bids are taken while it burns. When the flame gutters, the last bid wins. A failed Insight roll may mean you bid on a lot of inferior grade – the auctioneer’s smile is not a promise.

\subsection*{Iron Quotas (Nine Kingdomss)}
Each kingdom has an annual iron quota owed to the Aeler viceroyalty. Quotas are paid in ingots stamped with the viceroy’s mark. Smuggling unstamped iron is a capital offence – the forges of the Brass Coast pay a premium for “unregistered” metal, but the Aeler breath‑wardens can smell the difference.

\subsection*{Lighthouse Charters (Sidhi) as Currency}
On the Brass Coast, a Sidhi lighthouse charter is accepted as legal tender for port fees, and a Kuvani remount master may honour it as proof of credit. The Fhara water‑clerks are more suspicious – they weigh the wax seal against a water‑share tablet.

\subsection*{Monsoon Trade Winds}
From spring to autumn, the monsoon winds blow from the west, carrying Taharkan coffee and Ayokhan spices to Dhaharan ports. In winter, the winds reverse, carrying Dhaharan iron and brass east to Sihai and beyond. A captain who misses the window must pay a “wind toll” to the harbour master – or wait six months.

% −−−−−−−−−−−−−−−−−−−− ELEPHANTS & WAR −−−−−−−−−−−−−−−−−−−−
\chapter{Elephants & Siege Engines}

War elephants are not mere beasts; they are living siege towers. A fully armoured elephant can carry a howdah with archers, smash through wooden gates, and trample infantry lines.

\subsection*{Elephant Stats (Tier III)}
\begin{itemize}
  ℹtem \textbf{Harm:} 5 (stomp or tusk), \textbf{Fatigue:} Body 5
  ℹtem \textbf{Traits:} \textit{Trample} – when moving through Close range, each enemy must test Body + Athletics (DV 3) or take Harm 2 (crushing). \textit{Fear of Fire} – if targeted by a fire effect, the elephant must test Resolve (DV 4) or panic (stampede, uncontrolled).
  ℹtem \textbf{Mahout Bond:} An elephant bonded to a mahout gains +1 die to resist fear. The mahout can spend a Boon to reroll the elephant’s failed Resolve test.
\end{itemize}

\subsection*{Elephant Talents (for PCs)}
\begin{itemize}
  ℹtem \textbf{Mahout’s Whistle (2 XP)} – Once per scene, call your elephant to your side as a Move action.
  ℹtem \textbf{Howdah Archer (4 XP)} – When firing from a moving elephant, ranged attacks suffer no Position penalty.
  ℹtem \textbf{Tuskbreaker (6 XP)} – Your elephant’s trample attack ignores 1 point of armor.
\end{itemize}

% −−−−−−−−−−−−−−−−−−−− THE NINTH SEAL & THE KALRANTHA AMULET −−−−−−−−−−−−−−−−−−−−
\chapter{The Ninth Seal and the Kalrantha Amulet}

\subsection*{The Demon in the Mine}
Deep beneath the Himadri foothills, sealed by Covenant wax and Aeler iron, lies a fragment of \textbf{Varash the Unmaker} – a being of cosmic indifference that fed on the civilisation that turned the Red Desert to glass. It does not rage. It whispers. It offers bargains.

⫕section*{Whispers of the Seal (d6)}
When the party approaches the Ninth Seal or a fragment of the amulet, roll or choose a whisper:
\begin{enumerate}
  ℹtem \textit{“The seal was placed by a coward. He is still alive. Bring me his name, and I will give you a key that opens any lock – once.”}
  ℹtem \textit{“You carry a secret. I can hear it. Shall I keep it, or shall I tell the one who matters most?”}
  ℹtem \textit{“The Hollow and I are cousins. I taught it to count. It taught me to wait. Count to nine. I will show you the way out – but the way out leads through your own grave.”}
  ℹtem \textit{“The Vent‑Prior knows why the seal cracks. Ask him why he never counts his breaths aloud. He sealed me because he was afraid of what I would ask him to remember.”}
  ℹtem \textit{“The Kalrantha amulet is a fragment of me. Bring it here. I will make you whole. I will give you certainty – the absolute knowledge that every choice you make is the right one. (The whisper laughs.) It always was, wasn’t it? You just didn’t know.”}
  ℹtem \textit{“You are already dreaming this. When you wake, you will not remember me. But I will remember you. I will be waiting in the space between your heartbeats.”}
\end{enumerate}

\subsection*{The Kalrantha Amulet}
A shard of green glass set in a brass frame. It is said to grant a wish – but the wish always has a hidden clause. The Veiled Sisters want it to control the demon. The Breath‑less want to destroy it. The Covenant wants to seal it again. And the rajas of the Nine Kingdomss want to use it to settle old feuds.

⫕section*{Mechanics (GM only)}
The amulet is not a combat item. It is a plot device. Possessing it advances a \textbf{Kalrantha’s Attention} timer [8]. When the timer fills, the demon manifests as a Tier IV adversary – not to fight, but to bargain. The price is always a memory, a name, or a promise.

\subsection*{Integration with the Way of Silk and Eastern Realms}
\begin{itemize}
  ℹtem \textbf{Taharka and Ayokha (Monsoon Trade):} The amulet is said to have been carried east by a spice caravan that never reached its destination. Its trail goes cold in the Ayokhan highlands. A monsoon captain may have seen its green glow in a storm.
  ℹtem \textbf{Sihai (Eastern Realms):} The Lethai‑thora archives contain a scroll that describes the amulet’s creation – written by a mad emperor who believed he could bind a demon to secure his Mandate. The scroll is sealed with nine silk cords. The ninth cord has never been cut.
  ℹtem \textbf{The Brass Coast (Sidhi diaspora):} A Sidhi freedman claims her grandmother escaped from an Ashaani slave ship carrying the amulet. She threw it into the sea. The sea threw it back. It is now in the vault of a lighthouse judge, who uses it to blackmail the Admiralty.
\end{itemize}

% −−−−−−−−−−−−−−−−−−−− EPILOGUE −−−−−−−−−−−−−−−−−−−−
\chapter*{Epilogue: The Ninth Lease}
\addcontentsline{toc}{chapter}{Epilogue}

The lease is renewed every nine years. The iron quotas are adjusted. The spice auctions close at midnight. The demon whispers, and the mountain listens.

I have written eight scrolls for this book. The ninth is a lease – a contract between the Aeler viceroy and the raja of Kalinga, signed in blood and stamped with a bell that has no clapper. It expires next spring. I do not know if it will be renewed.

The road is south. The ink is fresh. The ninth seal is still warm.

\vspace{1em}
\noindent\textit{— Viceroy’s Scribe Govind, at the weigh‑house in Sarvistan, under a lamp that has just flickered for the eighth time}

\vspace{2em}
\begin{quote}
  \small
  \textit{“The elephant remembers. The mountain forgets. The demon waits. Choose your lease carefully.”} \\
  — Dusana of the Raven Road (marginal note, found in a Pereshi fire‑temple)
\end{quote}

% −−−−−−−−−−−−−−−−−−−−−−−−−−−−−−−−−−−−−−−−−−−−−−−−−−−−−−−−−−−−−−−−−−−−−−
% COLOPHON
% −−−−−−−−−−−−−−−−−−−−−−−−−−−−−−−−−−−−−−−−−−−−−−−−−−−−−−−−−−−−−−−−−−−−−−
\chapter*{Colophon}
\addcontentsline{toc}{chapter}{Colophon}

This book was written on paper that had been a cargo manifest, a love letter, and a bandage. The ink is lampblack cut with salt water and the ash of a burnt confession. The binding is grey cord tied with nine knots. The ninth knot is not cut.

Govind composed the preface in a weigh‑house while an Aeler breath‑warden counted his exhalations. The almanacs were gathered from merchants, miners, and the occasional honest raja. The generator tables were tested in the field – meaning he ran them while a spice auctioneer took bids.

If you find yourself in Dhahara and a viceroy’s scribe offers you a stamped receipt, accept it. If they offer you a lease, read the fine print. The ninth clause is always the one that kills you.

The elephant remembers. The mountain forgets. The demon waits.

\end{document}
